\worksheettitle
  {Отделяем классы справа}
  {\modulemark{M03-R} Коррекционная карточка}

\studentnameblock

\begin{notebox}[Как выполнять]
  Закройте число слева. Сначала выделите последние три цифры, затем следующие
  три. Левая группа может быть короче.
\end{notebox}

\section{Шаг 1. Отделяйте справа}

Разделите число \(8305072\).

\begin{enumerate}[label=\textbf{\arabic*.}]
  \item Последние три цифры: \answerblank[22mm].
  \item Следующие три цифры: \answerblank[22mm].
  \item Оставшаяся слева группа: \answerblank[18mm].
\end{enumerate}

Итог:
\[
  8305072=\answerblank[50mm].
\]

\section{Шаг 2. Сохраняйте три позиции}

В числе \(12\,005\,071\):

\begin{center}
\renewcommand{\arraystretch}{1.45}
\begin{tabular}{c|ccc}
  \toprule
  класс & миллионы & тысячи & единицы \\
  \midrule
  группа & 12 & \answerblank[16mm] & \answerblank[16mm] \\
  число класса & 12 & \answerblank[16mm] & \answerblank[16mm] \\
  \bottomrule
\end{tabular}
\end{center}

\begin{checkpointbox}[Повторная проверка]
  Разделите на классы: \(90305008=\answerblank[50mm].\)

  Класс единиц записан группой \answerblank[18mm] и содержит число
  \answerblank[18mm].
\end{checkpointbox}

